\chapter{Introdução}
\label{cap:introducao}

\section{Contextualização}
\label{sec:contextualizacao}

A formação em Engenharia de Controle e Automação exige a articulação entre
fundamentos teóricos --- sistemas de controle, eletrônica, programação --- e a
operação de sistemas físicos reais, como controladores lógicos programáveis
(CLPs), sensores, atuadores e painéis elétricos. Essa articulação
teórico-prática é, historicamente, um dos pontos críticos da formação: bancadas
didáticas completas têm custo elevado, exigem manutenção, ocupam espaço físico
e impõem riscos elétricos e mecânicos que restringem a experimentação livre
pelo estudante.

Paralelamente, os jogos digitais consolidaram-se como uma mídia de
aprendizagem. A literatura sobre jogos sérios demonstra que jogos podem unir
propósito educacional e engajamento \cite{abt1970serious}, que bons jogos
incorporam princípios de aprendizagem ativa que a escola persegue
\cite{gee2003what} e que a simulação interativa é uma ponte natural entre
treinamento, entretenimento e educação \cite{zyda2005from,
laamarti2014overview}. Um laboratório virtual não substitui a bancada real,
mas pode ocupar o espaço entre a teoria e a prática: oferecer familiaridade
com nomenclatura, procedimentos, intertravamentos e diagnóstico de falhas
antes --- e depois --- do contato com o equipamento físico.

Nos últimos anos, uma terceira transformação se somou às duas primeiras: os
grandes modelos de linguagem (LLMs) tornaram-se capazes de manter diálogos
naturais e de assumir papéis relevantes em jogos, da geração de conteúdo à
animação de personagens \cite{gallotta2024largelanguage}. Essa capacidade abre
uma possibilidade didática específica: um personagem de IA que comenta o
estado do laboratório, responde a perguntas técnicas e reage às decisões do
jogador --- um instrutor ficcional sempre disponível. Ao mesmo tempo, ela
traz dois riscos práticos. O primeiro é de \emph{controle}: como garantir que
a IA enriqueça o jogo sem, em nenhum momento, decidir ou executar ações?
Experiências conhecidas com personagens generativos e agentes autônomos
mostram que muita autonomia produz comportamentos ricos, mas imprevisíveis
\cite{park2023generative,wang2023voyager}. O segundo é de
\emph{infraestrutura}: modelos grandes de linguagem normalmente exigem
serviços externos ou hardware especializado, o que contraria a realidade de
laboratórios escolares modestos.

Este trabalho nasce exatamente nesse encontro. Ele descreve o projeto, o
desenvolvimento e a validação técnica do \textbf{Laboratório 404}, um jogo
sério 3D educacional/ficcional no qual uma IA local, a personagem \aria{},
participa do mundo por meio de diálogo, sem nunca controlar o jogo. O
protótipo foi construído de forma incremental, em seis provas de conceito
(\poc{1} a \poc{6}), sobre ferramentas abertas --- o motor Godot, o Blender para
os assets 3D e o Ollama para a inferência local --- e avaliado em um
computador pessoal sem GPU dedicada.

\section{Problema}
\label{sec:problema}

O problema que orienta este trabalho pode ser enunciado em três dimensões
interligadas:

\begin{enumerate}
  \item \textbf{Segurança e previsibilidade da IA.} Um jogo educacional
  precisa que cada ação do sistema seja reproduzível e auditável. É necessário
  integrar um LLM a um jogo 3D de modo que ele enriqueça o mundo e a narrativa
  sem receber qualquer capacidade de alterar o estado do jogo.
  \item \textbf{Viabilidade em hardware modesto.} A integração precisa
  funcionar em um computador pessoal sem GPU dedicada, mantendo o desempenho
  gráfico adequado (jogabilidade fluida) e latências de conversa toleráveis
  para o tipo de interação proposto.
  \item \textbf{Rastreabilidade do processo.} O desenvolvimento --- de autor
  único, assistido por agentes de inteligência artificial --- precisa ser
  conduzido e documentado de forma que cada afirmação sobre o protótipo seja
  sustentada por requisito, teste ou evidência verificável.
\end{enumerate}

Essas dimensões se desdobram nas questões de pesquisa que este trabalho
persegue: (RQ1) como arquitetar a integração de um LLM local a um jogo 3D de
modo que ele participe do mundo sem controlá-lo? (RQ2) é possível manter
jogabilidade fluida e conversa local com latência tolerável em um computador
sem GPU dedicada? (RQ3) como conduzir e documentar esse desenvolvimento de
forma rastreável e reprodutível?

\section{Justificativa}
\label{sec:justificativa}

A justificativa acadêmica deste trabalho repousa em três argumentos. O
primeiro é a lacuna de integração: a literatura concentra-se, de um lado, em
agentes autônomos capazes de agir \cite{park2023generative,wang2023voyager} e,
de outro, em visões amplas sobre o papel das LLMs em jogos
\cite{gallotta2024largelanguage}; é menos explorado o desenho oposto, de uma
IA diegética com capacidade formalmente limitada, integrada a um jogo
educacional completo e documentada com evidências de validação.

O segundo argumento é a relevância prática para o ensino de automação: um
laboratório simulado de baixo custo, executado em hardware comum e sem
dependência de serviços pagos, amplia o tempo de exposição do estudante a
situações de operação, falha e diagnóstico --- situações que, na bancada
física, são limitadas por segurança, custo e disponibilidade.

O terceiro argumento é metodológico: o trabalho documenta detalhadamente um
processo de desenvolvimento orientado a especificação, com critérios de
aceitação verificáveis, suíte automatizada e evidências visuais, incluindo as
limitações e as pendências reais do protótipo. Esse registro tem valor para
outros projetos que precisem integrar IA generativa a sistemas interativos com
previsibilidade.

\section{Objetivos}
\label{sec:objetivos}

\subsection{Objetivo geral}

Desenvolver e validar tecnicamente um jogo sério 3D, construído sobre
tecnologias abertas e executável em computador pessoal sem GPU dedicada, no
qual uma inteligência artificial local participe do mundo por meio de diálogo
--- sem controlar o jogo --- e que apoie, de forma ficcional e procedural, o
ensino de conceitos de automação industrial.

\subsection{Objetivos específicos}

\begin{enumerate}
  \item Especificar o sistema com requisitos funcionais e não funcionais,
  critérios de aceitação no formato DADO/QUANDO/ENTÃO e matriz de
  rastreabilidade requisito--teste;
  \item estruturar o desenvolvimento em provas de conceito incrementais, cada
  uma terminando em versão executável e demonstrável;
  \item implementar o pipeline de assets 3D de modelagem no Blender até a
  importação no motor Godot, com convenções de escala, nomenclatura e colisão;
  \item implementar a integração com a IA local em processo separado, com
  contrato validado, intenções restritas e resposta de \emph{fallback};
  \item implementar o mundo reativo (estados de equipamento, missão, NPC,
  memória, iluminação) e a geração procedural determinística da ala final;
  \item construir a suíte de validação automatizada e o arnês de evidências
  visuais, executá-los e registrar os resultados;
  \item medir o desempenho gráfico e as latências da conversa na máquina de
  referência, e registrar as limitações e pendências do protótipo.
\end{enumerate}

\section{Delimitação}
\label{sec:delimitacao}

Este trabalho é um estudo de caso de desenvolvimento tecnológico. Ele
\emph{não} avalia aprendizagem nem experiência de usuários: não há
experimento educacional controlado, nem estudo de usabilidade com
participantes. A validação realizada é funcional e técnica: requisitos,
testes automatizados, medições de desempenho e latência e evidências visuais.
O escopo do protótipo é um recorte jogável de 10 a 20 minutos, com uma missão
principal, uma ala procedural e um epílogo de escolhas, executado em uma
única máquina de referência com GPU integrada. A validação em hardware real
(joystick e mouse) e o empacotamento independente permanecem pendentes e são
declarados como tal.

\section{Visão geral da metodologia}
\label{sec:metodologia_visao}

O trabalho adota uma abordagem aplicada e incremental. A especificação foi
escrita antes da implementação e funcionou como contrato: requisitos,
critérios de aceitação e regras de projeto foram fixados em documentos
versionados, e cada iteração de desenvolvimento precisava satisfazer um
subconjunto verificável desses critérios. Seis provas de conceito organizaram
a construção: fundação jogável (\poc{1}); laboratório funcional com missão
(\poc{2}); integração da IA local (\poc{3}); mundo reativo (\poc{4});
laboratório procedural (\poc{5}); e recorte vertical final (\poc{6}). A cada
etapa, a suíte automatizada \emph{headless} e um conjunto de testes do
adaptador de IA foram executados, e evidências visuais foram capturadas por um
arnês de captura de tela com posicionamento programático do jogador. O
detalhamento do método está no \autoref{cap:metodos}.

\section{Organização do trabalho}
\label{sec:organizacao}

Além desta introdução, o trabalho está organizado em sete capítulos. O
\autoref{cap:fundamentacao} apresenta a fundamentação teórica; o
\autoref{cap:relacionados} discute os trabalhos relacionados; o
\autoref{cap:metodos} descreve materiais e métodos; o
\autoref{cap:desenvolvimento} detalha o desenvolvimento do sistema; o
\autoref{cap:resultados} apresenta os experimentos e resultados; o
\autoref{cap:discussao} discute os resultados e as limitações; e o
\autoref{cap:conclusao} conclui o trabalho e aponta trabalhos futuros. Os
apêndices reúnem os requisitos completos, o detalhamento da suíte de testes e
a descrição do repositório.
